\section{行省之内：地方文书、河运与晚元秩序}

\subsection{行省不是地图上的色块}

元代行中书省处理的是一组不断移动的问题：战后怎样接收旧有州县，军粮怎样抵达前线，赋税怎样送进仓场，案件怎样由地方送到能够裁决的官员手中。它并非今天意义上的固定省份，也不能从一张制度图推知每处地方的实际行政强度。行省官员、路府州县、军政机构、站赤与地方中介之间的权限会因战争、河道、皇帝的调动和人事关系反复重叠。一个机构的设立只能说明朝廷试图建立接口，不能证明接口已经把所有村落接入同一种秩序。

江南接收以后，这个问题尤其清楚。新朝需要南宋留下的田赋、盐课、仓库和书写人员，却不能把它们直接装进大都的国库。征粮要依靠熟悉乡里的里正、书手和仓吏，运输要经过河港、堤岸、船户和转运节点；任何一处延误都可能使“应纳”与“实收”成为两种数目。地方官在奏报中会以欠额、灾伤或逃户说明困难，中央则用限期、复核和处分要求责任。这些文字让我们看见国家怎样追问资源，却不能完整复原一户人家如何在收成、租佃、借贷和劳役之间安排生计。

华北、江南、云南和东北的差别也不能被“行省制”一词抹平。华北的战后荒地、军镇和通往大都的道路，江南的水网、市场与旧有赋役传统，西南的山地首领和长距离交通，东北的边境与高丽联系，各自使文书、军队和征发面对不同阻力。朝廷可以使用相近的官名，却未必拥有相近的信息；同一条条格在都城附近、河谷城镇和山地聚落之间，也会经过不同的翻译、解释与协商。多区域治理的能力，正由这种不均匀而非由名目的一致来衡量。

\subsection{文书里的中介人}

《元典章》《通制条格》使人容易先看到皇帝、行省和法律的语言；但一份命令要变成地方行动，还要经过抄写、传递、释读、登记、核验和回报。书吏、译者、通事、里正、仓吏、牙人和担保人很少像大臣一样进入本纪，却可能决定一项税粮是否被正确登记、一件诉讼是否能进入衙门、一名使者是否能在驿路上获得马匹。将这些人简单称为腐败的阻碍或国家的工具，都不足以说明他们在信息不对称中拥有的实际选择。

多语环境使这种中介更加重要。蒙古文、汉文、畏兀儿文、八思巴字、藏文和波斯文分别在不同机构、宗教网络与贸易场合留下痕迹。文字并不是稳定身份的标签：同一个人可能因官署、市场或寺院而使用不同文字，官号的转写也未必能说明实际职权。碑刻、印章和文书能够校正正史中被汉译或简化的名称，却各自有保存地点和文类的限制。它们提示的是帝国如何依赖翻译与书写，而不是一种自然和谐的多民族行政。

司法文书也应这样读。条格规定什么可诉、何种证据可采，显示的是制度希望怎样分类债务、田宅、婚姻、逃役和商业；具体案件才可能让我们看见某次争执如何被说成合法语言。可是进入案件的人已经拥有一定的交通、文书和关系条件，未能进入衙门的冲突难以留下记录。不能以一条法律推论所有人受同样保护，也不能以一件冤案判定整个法制无效。真正值得追问的是，谁有能力把自己的经历变成一份可被官府读取的文本。

\subsection{纸钞、漕粮与海运的三种时间}

纸钞在元朝财政中提供了跨区域支付的可能，但纸上的数额、市场的价格和仓场里的实物并不处在同一种时间。朝廷以钞赏赐或规定折纳，可以迅速改变账簿；商人接受某种票钞、农户能否以所得换回生活所需、军队能否按时得到粮饷，却取决于流通、兑换和地方信任。钞法反复调整本身说明国家在处理这种距离，不能自动证明一项新规定在所有市镇获得相同效果。

漕粮又有不同的节奏。江南的粮食必须经过收割、征收、仓储、装船、过闸、转运和北方发放，才会成为大都的供给。河流的水位、堤岸的损坏、盗抢、风暴和船只的调配，会让一年的税额在不同地点变成截然不同的生活压力。海运在某些时期提供替代路线，却同样依赖港口、船材、航海者和沿岸秩序。把河运或海运写成技术胜利，会忽略每一次运输都要占用船户、工匠、仓夫和地方市场的资源。

盐课、商税和市舶收入也不能从财政表格直接化约为国家强盛。它们来自具体商品、路线和征收关系，受到市场需求、海上风险、地方保护与官员执行的影响。港口碑刻、沉船和钱币能让部分货物与社群进入视野，不能计算帝国全部贸易额；奏议中的收入忧虑则反映中枢怎样看待可征之财，不能替商人、船户和消费者说明成本。财政史在这里不是统计数字的竞赛，而是追问一项收入在何处、由谁、以什么代价才变得可被朝廷看见。

\subsection{宗教机构与地方资源}

佛寺、道观、伊斯兰礼拜空间、藏传佛教机构和地方祠祀，在元代既是仪式场所，也可能连接土地、施主、工匠、旅人、刻书与慈善。国家对僧道的保护、登记和限制，说明统治者希望这些机构如何进入资源与秩序；它们不能证明所有寺院受同样控制，更不能把某一位皇帝的宗教倾向直接等同于地方信仰。寺院题记常记录重修、捐施和愿望，因而更容易保存有财力留名者的声音。

宗教网络在交通不稳的时期尤其具有实际意义。旅人可能依靠熟悉的礼拜或住宿关系寻找信息，地方家族可能通过捐施建立纪念与互助，书写者也可能在经卷、碑石和法会中留下语言的痕迹。这不意味着寺院是脱离国家的避难所：赋役、土地、军队和官府会持续介入其资源。更准确的说法是，宗教机构与行政、市场和家族并行运作，并在不同地区形成不同的协商边界。

元朝的多语宗教世界也提醒我们，不应把“外来”与“本地”写成两种静止实体。伊斯兰、佛教、道教、景教和地方信仰的实践，既可能沿商路和迁徙传播，也会在具体城镇、婚姻、工艺和施主关系中重新扎根。材料能确认的往往是一座建筑、一通题记或一份文书，不能由此推论整个地区的信徒比例或彼此关系。承认这种尺度限制，并非放弃解释，而是防止帝国的丰富性被一串宗教名称取代。

\subsection{晚元：地方秩序怎样失去共同的节拍}

十四世纪中叶的危机并非从一个原因开始。黄河决溢和治河动员增加了北方村落、渡口与漕运的压力；纸钞、税粮和军费的困难使中枢更急于调动资源；宫廷人事更替又改变了谁能协调行省、军队和仓储。红巾诸集团在不同地区兴起时，所利用的宗教语言、地方武装、粮价和交通条件并不相同。把这些复杂情形压缩成“河工导致起事”或“财政崩溃导致灭亡”，都把同时发生误作单一因果。

脱脱被罢之后，军政调度的接口并没有随之自动恢复。地方将领、行省官员和区域势力必须在粮道、城防、征兵与市场之间寻找自己的资源，江南、水陆交通线和北方军镇的变化也不再由一个中枢节奏统一。正史擅长记录罢免、战役和城市易手，却难以呈现普通家庭怎样在逃避征发、迁居、复耕或加入新的保护网络之间选择。正因为如此，晚元社会不能只作为下一王朝胜利的背景。

一三六八年大都失守，是元朝中原统治的明确终点，却不是所有区域在同一天改变生活的终点。北方草原、东北、高丽关系、江南城市和海上网络各有自己的后续时间。新的王朝会重新登记、征税和安排道路，但旧制度留下的行省经验、运输线路、多语书写和地方关系并不会立即消失。元朝的终局由此也是一个观察制度遗产的时刻：它所解决过的跨区域协调问题，并未因一座都城易手而从中国和欧亚的历史中消退。
